\begin{dang}{Bài tập cơ bản}
\Anlythuyet{	\Binhluan{
	\begin{itemize}
	\item Khi máy phát có số cặp cực thay đổi $\Delta p$ và số vòng quay thay đổi $\Delta n$ (nên đổi đơn vị là vòng/giây) thì tùy thuộc trường hợp để lựa chọn dấu ''+'' hay dấu ''$-$'' trong các công thức sau:
	\begin{align*} 
	&\Delta n\,\left( {{\textrm{vòng/s}}} \right)\\
	&{f_1} = {n_1}{p_1} \Rightarrow {n_1} = \dfrac{{{f_1}}}{{{p_1}}}\\
	&{f_2} = {n_2}{p_2} = \left( {{n_1} \pm \Delta n} \right)\left( {{p_1} \pm \Delta p} \right) \Rightarrow {p_1} =?
	\end{align*}
	\item  Tổng số vòng dây của phần ứng $N = \dfrac{E_0}{\omega \Phi _0}$. 
\end{itemize}
}
{
Nếu phần ứng gồm k cuộn dây giống nhau mắc nối tiếp thì số vòng dây trong mỗi cuộn: $N_1 = \dfrac{N}{k}$.
}}
\end{dang}
